\thispagestyle{fancy}
\fancyhead[R]{A.U 2025-2026}
\fancyhead[L]{Annexes}

% boîte utilisée pour encadrer les extraits de code (verbatim)
\newsavebox{\annexbox}

\chapter{Installation et configuration de l'agent}

L'agent de collecte est un programme Python autonome à déployer sur chaque serveur à
superviser. Son installation se déroule en quatre étapes.

\section{Prérequis}
\begin{itemize}[label=\ding{43}]
  \item Python 3.8 ou version ultérieure.
  \item Un accès réseau sortant vers le backend de la plateforme.
  \item Les droits nécessaires à la création d'un service \texttt{systemd}.
\end{itemize}

\section{Installation des dépendances}

\begin{lrbox}{\annexbox}
\begin{minipage}{0.88\textwidth}
\begin{verbatim}
cd /opt/monitoring-agent
python3 -m venv venv
source venv/bin/activate
pip install -r requirements.txt
\end{verbatim}
\end{minipage}
\end{lrbox}
\begin{figure}[H]
\centering
\fbox{\usebox{\annexbox}}
\caption{Installation des dépendances de l'agent}
\label{fig:annexe-install}
\end{figure}

\section{Fichier de configuration}
Le comportement de l'agent est entièrement piloté par le fichier \texttt{config.json},
présenté dans la figure~\ref{fig:annexe-config}.

\begin{lrbox}{\annexbox}
\begin{minipage}{0.88\textwidth}
\begin{verbatim}
{
  "server_id": "srv-prod-01",
  "server_name": "Serveur de production",
  "backend_url": "http://141.227.129.194:30300/metrics",
  "interval": 5,
  "timeout": 10,
  "detect_services": true
}
\end{verbatim}
\end{minipage}
\end{lrbox}
\begin{figure}[H]
\centering
\fbox{\usebox{\annexbox}}
\caption{Fichier de configuration \texttt{config.json} de l'agent}
\label{fig:annexe-config}
\end{figure}

Le tableau~\ref{tab:annexe-config} détaille la signification de chaque paramètre.

\renewcommand{\arraystretch}{1.3}
\begin{longtable}{|>{\raggedright\arraybackslash}p{3.6cm}
                  |>{\raggedright\arraybackslash}p{10.2cm}|}
\hline
\textbf{Paramètre} & \textbf{Description} \\
\hline
\endfirsthead

\hline
\textbf{Paramètre} & \textbf{Description} \\
\hline
\endhead

\hline
\caption{Paramètres de configuration de l'agent}
\label{tab:annexe-config}
\endlastfoot

\texttt{server\_id} & Identifiant unique du serveur, utilisé comme clé côté backend. \\
\hline
\texttt{server\_name} & Libellé affiché dans l'interface d'administration. \\
\hline
\texttt{backend\_url} & Adresse de l'endpoint de réception des métriques. \\
\hline
\texttt{interval} & Intervalle de collecte, exprimé en secondes. \\
\hline
\texttt{timeout} & Délai maximal d'attente d'une réponse du backend, en secondes. \\
\hline
\texttt{detect\_services} & Active ou désactive la détection automatique des services actifs. \\
\hline

\end{longtable}

\section{Déclaration du service systemd}
Afin que l'agent démarre automatiquement avec le système et redémarre en cas d'arrêt
inopiné, il est déclaré comme service \texttt{systemd}.

\begin{lrbox}{\annexbox}
\begin{minipage}{0.88\textwidth}
\begin{verbatim}
[Unit]
Description=Agent de supervision CLEDISS
After=network.target

[Service]
Type=simple
User=monitoring
WorkingDirectory=/opt/monitoring-agent
ExecStart=/opt/monitoring-agent/venv/bin/python3 main.py
Restart=always
RestartSec=10

[Install]
WantedBy=multi-user.target
\end{verbatim}
\end{minipage}
\end{lrbox}
\begin{figure}[H]
\centering
\fbox{\usebox{\annexbox}}
\caption{Unité \texttt{systemd} de l'agent de supervision}
\label{fig:annexe-systemd}
\end{figure}

L'activation du service s'effectue ensuite par les commandes suivantes.

\begin{lrbox}{\annexbox}
\begin{minipage}{0.88\textwidth}
\begin{verbatim}
sudo systemctl daemon-reload
sudo systemctl enable monitoring-agent
sudo systemctl start monitoring-agent
sudo systemctl status monitoring-agent
\end{verbatim}
\end{minipage}
\end{lrbox}
\begin{figure}[H]
\centering
\fbox{\usebox{\annexbox}}
\caption{Activation du service de l'agent}
\label{fig:annexe-activation}
\end{figure}

\chapter{Configuration du backend}

\section{Variables d'environnement}
La configuration du backend est intégralement externalisée sous forme de variables
d'environnement, injectées en production depuis le coffre HashiCorp Vault. Le
tableau~\ref{tab:annexe-env} en présente la liste.

\renewcommand{\arraystretch}{1.3}
\begin{longtable}{|>{\raggedright\arraybackslash}p{4.2cm}
                  |>{\raggedright\arraybackslash}p{9.6cm}|}
\hline
\textbf{Variable} & \textbf{Description} \\
\hline
\endfirsthead

\hline
\textbf{Variable} & \textbf{Description} \\
\hline
\endhead

\hline
\caption{Variables d'environnement du backend}
\label{tab:annexe-env}
\endlastfoot

\texttt{PORT} & Port d'écoute du serveur Express (3000 par défaut). \\
\hline
\texttt{MONGODB\_URI} & Chaîne de connexion authentifiée à la base MongoDB Atlas. \\
\hline
\texttt{JWT\_SECRET} & Clé de signature et de vérification des jetons d'authentification. \\
\hline
\texttt{EMAIL\_USER} & Adresse du compte SMTP émetteur des notifications. \\
\hline
\texttt{EMAIL\_PASS} & Mot de passe d'application du compte SMTP. \\
\hline
\texttt{ADMIN\_EMAIL} & Adresse du compte administrateur créé à l'initialisation. \\
\hline
\texttt{ADMIN\_PASSWORD} & Mot de passe du compte administrateur initial. \\
\hline

\end{longtable}

\section{Seuils de déclenchement des alertes}
Les seuils sont stockés en base de données et modifiables sans redéploiement. Leurs
valeurs par défaut figurent dans le tableau~\ref{tab:annexe-seuils}.

\renewcommand{\arraystretch}{1.3}
\begin{longtable}{|>{\raggedright\arraybackslash}p{4.6cm}
                  |>{\centering\arraybackslash}p{3.4cm}
                  |>{\centering\arraybackslash}p{3.4cm}|}
\hline
\textbf{Métrique} & \textbf{Seuil \texttt{WARNING}} & \textbf{Seuil \texttt{CRITICAL}} \\
\hline
\endfirsthead

\hline
\textbf{Métrique} & \textbf{Seuil \texttt{WARNING}} & \textbf{Seuil \texttt{CRITICAL}} \\
\hline
\endhead

\hline
\caption{Seuils de déclenchement des alertes}
\label{tab:annexe-seuils}
\endlastfoot

Utilisation du processeur & 80 \% & 90 \% \\
\hline
Utilisation de la mémoire vive & 80 \% & 90 \% \\
\hline
Utilisation du disque & 80 \% & 90 \% \\
\hline
Absence de métrique reçue & 2 minutes & 5 minutes \\
\hline

\end{longtable}

\chapter{Commandes d'exploitation}

\section{Cluster Kubernetes}

\begin{lrbox}{\annexbox}
\begin{minipage}{0.88\textwidth}
\begin{verbatim}
# État général du cluster
kubectl get nodes
kubectl get all -n monitoring

# Journaux d'un pod
kubectl logs -f deployment/backend -n monitoring

# Redémarrage progressif d'un déploiement
kubectl rollout restart deployment/backend -n monitoring
kubectl rollout status deployment/backend -n monitoring

# Retour à la version précédente
kubectl rollout undo deployment/backend -n monitoring
\end{verbatim}
\end{minipage}
\end{lrbox}
\begin{figure}[H]
\centering
\fbox{\usebox{\annexbox}}
\caption{Commandes d'exploitation du cluster k3s}
\label{fig:annexe-kubectl}
\end{figure}

\section{Conteneurs Docker}

\begin{lrbox}{\annexbox}
\begin{minipage}{0.88\textwidth}
\begin{verbatim}
# Construction des images
docker build -t monitoring-backend:latest ./backend
docker build -t monitoring-frontend:latest ./frontend
docker build -t monitoring-agent:latest ./agent

# Démarrage de la pile complète en local
docker compose up -d
docker compose logs -f backend
docker compose down

# Analyse de vulnerabilites d'une image
trivy image --severity HIGH,CRITICAL monitoring-backend:latest
\end{verbatim}
\end{minipage}
\end{lrbox}
\begin{figure}[H]
\centering
\fbox{\usebox{\annexbox}}
\caption{Commandes de construction et d'analyse des images}
\label{fig:annexe-docker}
\end{figure}

\chapter{Récapitulatif des services web}

Le tableau~\ref{tab:annexe-api} récapitule l'ensemble des services web exposés par la
plateforme à l'issue des quatre sprints de réalisation.

\renewcommand{\arraystretch}{1.3}
\small
\begin{longtable}{|>{\raggedright\arraybackslash}p{1.6cm}
                  |>{\raggedright\arraybackslash}p{7.4cm}
                  |>{\centering\arraybackslash}p{2.4cm}
                  |>{\centering\arraybackslash}p{2.6cm}|}
\hline
\textbf{Méthode} & \textbf{URL} & \textbf{Sprint} & \textbf{Auth.} \\
\hline
\endfirsthead

\hline
\textbf{Méthode} & \textbf{URL} & \textbf{Sprint} & \textbf{Auth.} \\
\hline
\endhead

\hline
\multicolumn{4}{r}{\textit{Suite du tableau à la page suivante}} \\
\endfoot

\hline
\caption{Récapitulatif de l'ensemble des services web de la plateforme}
\label{tab:annexe-api}
\endlastfoot

\texttt{POST} & \texttt{/api/auth/login} & 1 & -- \\
\hline
\texttt{GET} & \texttt{/api/auth/me} & 1 & JWT \\
\hline
\texttt{POST} & \texttt{/metrics} & 1 & -- \\
\hline
\texttt{GET} & \texttt{/api/servers} & 1 & JWT \\
\hline
\texttt{POST} & \texttt{/api/servers} & 1 & JWT \\
\hline
\texttt{GET} & \texttt{/api/servers/:server\_id} & 1 & JWT \\
\hline
\texttt{PUT} & \texttt{/api/servers/:server\_id} & 1 & JWT \\
\hline
\texttt{GET} & \texttt{/api/servers/:server\_id/metrics} & 1 & JWT \\
\hline
\texttt{GET} & \texttt{/api/metrics/latest} & 1 & JWT \\
\hline
\texttt{GET} & \texttt{/api/metrics/history/:serverId} & 1 & JWT \\
\hline
\texttt{GET} & \texttt{/api/metrics/stats} & 1 & JWT \\
\hline
\texttt{GET} & \texttt{/api/dashboard/summary} & 1 & JWT \\
\hline
\texttt{GET} & \texttt{/api/alerts} & 2 & JWT \\
\hline
\texttt{GET} & \texttt{/api/alerts/:alert\_id} & 2 & JWT \\
\hline
\texttt{PUT} & \texttt{/api/alerts/:alert\_id/acknowledge} & 2 & JWT \\
\hline
\texttt{PUT} & \texttt{/api/alerts/:alert\_id/resolve} & 2 & JWT \\
\hline
\texttt{GET} & \texttt{/api/alerts/stats/summary} & 2 & JWT \\
\hline
\texttt{POST} & \texttt{/api/alerts/bulk/acknowledge} & 2 & JWT \\
\hline
\texttt{GET} & \texttt{/api/remote-actions/}\allowbreak\texttt{:server\_id/}\allowbreak\texttt{services-status} & 2 & JWT + admin \\
\hline
\texttt{POST} & \texttt{/api/remote-actions/}\allowbreak\texttt{:server\_id/}\allowbreak\texttt{start-service} & 2 & JWT + admin \\
\hline
\texttt{POST} & \texttt{/api/remote-actions/}\allowbreak\texttt{:server\_id/}\allowbreak\texttt{stop-service} & 2 & JWT + admin \\
\hline
\texttt{POST} & \texttt{/api/remote-actions/}\allowbreak\texttt{:server\_id/}\allowbreak\texttt{restart-service} & 2 & JWT + admin \\
\hline
\texttt{POST} & \texttt{/api/remote-actions/}\allowbreak\texttt{:server\_id/}\allowbreak\texttt{restart} & 2 & JWT + admin \\
\hline
\texttt{POST} & \texttt{/api/remote-actions/}\allowbreak\texttt{:server\_id/}\allowbreak\texttt{shutdown} & 2 & JWT + admin \\
\hline
\texttt{GET} & \texttt{/api/remote-actions/}\allowbreak\texttt{:server\_id/}\allowbreak\texttt{audit-log} & 2 & JWT + admin \\
\hline
\texttt{GET} & \texttt{/api/backups} & 2 & JWT \\
\hline
\texttt{GET} & \texttt{/api/backups/:id} & 2 & JWT \\
\hline

\end{longtable}
\normalsize
